% ============================================================
\section{Related Work}
\label{sec:related}
% ============================================================

\noindent\textbf{Post-training quantization and bit allocation.}
PTQ has matured from uniform rounding \cite{gptq,adaround,repqvit} and
learned step sizes \cite{jacob2018quant,esser2020lsq}
to mixed precision, where layers receive different bit-widths under a size or
bit budget \cite{haq,hawq,bitsplit,dong2020hawqv2,yao2021hawqv3},
including reconstruction-based \cite{nagel2021brecq,wei2022qdrop}
and activation-aware LLM objectives \cite{xiao2023smoothquant,lin2024awq}.
Most mixed-precision methods follow
the score-then-allocate template, which itself inherits the classical
sensitivity-pruning lineage \cite{lecun1990obd}: compute
a per-layer sensitivity---often
Hessian- or Fisher-based \cite{hawq,lampq}, sometimes
task-metric- or
regression-based \cite{regptq}---and then solve a budgeted program (DP,
ILP, or Lagrangian search) over these independent scores.

\smallskip
\noindent\textbf{Detector-specific quantization.}
Detector backbones have evolved from CNN detectors
\cite{ren2015faster} to transformer detectors
\cite{carion2020detr,zhu2021deformable} and ViT-style backbones
\cite{dosovitskiy2021vit}, with dedicated PTQ designs
\cite{li2022ptq4vit,liu2021qvit,lin2022fqvit}.
F-LBQ \cite{flbq} allocates transformer-layer bit-widths by minimizing
backbone-feature alteration (uniform quantization on heads; layer-wise
feature-space proxy). LampQ \cite{lampq} uses a type-aware Fisher-style
score with an ILP allocation and reports Mask R-CNN + Swin-T detection
results at 4-bit weight/activation. Reg-PTQ \cite{regptq} targets the
regression branch with a global reconstruction loss. These works share
the singleton-score premise; we instead measure \emph{conditional}
utility and show the singleton premise is quantitatively unreliable for
detector allocation decisions.

\smallskip
\noindent\textbf{Output-distortion objectives.}
BLOB-Q \cite{blobq} approximates global distortion by a second-order
expansion and, under a layer-uncorrelatedness assumption, decomposes it
into additive terms enabling DP allocation; it reports Mask R-CNN +
Swin-T/S detection at 4/5/6-bit. We test the
\emph{operational consequence} of additive decomposition: can
context-free contributions predict which block deserves extra precision
\emph{after} others changed? Our interaction and rank-reversal
measurements indicate they cannot (Sections~\ref{sec:mech},
\ref{sec:exp}), which motivates measuring marginals in context.

\smallskip
\noindent\textbf{Calibration without labels.}
Most quantization calibration uses unlabeled data
\cite{cai2020zeroq,krishnamoorthi2018}; what varies is the
quantity measured. Feature-space proxies \cite{flbq}, output-level
quadratic objectives \cite{blobq}, and gradient/Fisher statistics
\cite{lampq} can all be computed without labels. Ours also needs no
labels---only FP-teacher detections---but measures \emph{task-output
agreement in the current context} rather than a singleton proxy, with a
sample-size analysis (Sec.~\ref{sec:theory}).

\smallskip
\noindent\textbf{Reimplementation note.} Prior-art comparators are
labeled reimplementations/stand-ins (feature-MSE for F-LBQ-style,
additive output sensitivity for BLOB-Q-style, and a Fisher-style score
for LampQ-style); they share our quantizer, calibration set and budget
accounting, and are not the official pipelines.
